\documentclass[11pt]{article}

\usepackage[margin=1in]{geometry}
\usepackage[utf8]{inputenc}
\usepackage[T1]{fontenc}
\usepackage{amsmath,amssymb}
\usepackage[colorlinks=true,linkcolor=blue,citecolor=blue,urlcolor=blue]{hyperref}
\usepackage{enumitem}
\usepackage{fancyhdr}

\pagestyle{fancy}
\fancyhf{}
\fancyfoot[C]{\small Preprint --- DOI: 10.5281/zenodo.22261524}
\fancyfoot[R]{\small \thepage}
\renewcommand{\headrulewidth}{0pt}

\title{\textbf{Selfhood Before Heredity: From a Peptide QAG to the Percolating Layer L*}}

\author{Alex Shoshitaishvili\\ \textit{Independent Researcher}}

\date{September 2026}

\begin{document}

\maketitle

\begin{center}
\textit{Preprint. DOI: \href{https://doi.org/10.5281/zenodo.22261524}{10.5281/zenodo.22261524}}\\
\textit{Repository: \url{https://github.com/alexshosh/selfhood-before-heredity}}
\end{center}

\vspace{1em}

\begin{abstract}
\noindent
The astronomical size of protein sequence space poses a needle-in-a-haystack problem: how could functional peptide and protein organization become accessible without exhaustive search? We seek a peptide/protein configuration-space layer, L*, in which viable configurations form a sufficiently connected network for sustained evolutionary traversal. The project approaches L* from two directions. From the modern side, the GB1 domain provides an experimental anchor: independently supplied interacting peptide fragments can reconstitute a native-like complex, providing seed systems for mapping viable connectivity {[}14{]}. From the geochemical side, Deacon's autogenic principle shows how reciprocally linked self-organizing processes can constitute a reproducing chemical self without prior sequence heredity {[}12,13{]}. Matsuo and Kurihara provide an experimental peptide-coacervate realization of this topology, which we classify as a quasi-autogenic organization (QAG) {[}1{]}. Although their implementation uses laboratory feed chemistry, experimentally supported cysteine, carbonyl-sulfide, thiol, thioester, and peptide-forming chemistry provides a reaction-space connection to geochemical matter {[}1-5,16,17{]}. We propose that these two approaches can enter L* at different points. If both points belong to the same percolating viable network, an evolutionary path can connect the geochemically grounded QAG side to the modern-protein side without requiring exhaustive search or prior sequence heredity. We also ground the synthetic QAG demonstrated by Matsuo and Kurihara {[}1{]} in geochemical reaction space.
\end{abstract}

\hypertarget{introduction-the-argument-in-outline}{%
\section*{Introduction: the argument in outline}\label{introduction-the-argument-in-outline}}

The enormous size of protein sequence space creates a needle-in-a-haystack problem: life could not have begun by randomly searching for something resembling a modern protein. The relevant alternative is percolation. If viable peptide configurations form a sufficiently connected network, evolution can move through that network by successive viable changes without searching the surrounding astronomical space. We therefore seek a peptide/protein configuration-space layer, L*, containing such a connected viable network. Importantly, the two approaches to L* need not meet at the same molecular configuration. From above, the modern-protein route need only reach a viable interacting peptide system within this connected region. From below, the geochemical route need only reach a sufficiently developed quasi-autogenic organization within that same region. Once both points belong to the same percolating viable network, a horizontal evolutionary path through configuration space can in principle connect them. Thus we do not need to construct a direct chemical or historical route from a primitive peptide system to a particular modern protein; we need to establish that the upward and downward constructions can enter the same connected evolutionary region.

Deacon supplies the organizational principle needed for the upward route: reciprocally linked self-organizing processes can constitute a reproducing chemical self without requiring prior sequence heredity. Matsuo and Kurihara provide an experimental peptide-coacervate realization of this principle, which we classify as a quasi-autogenic organization (QAG). Their particular laboratory implementation is not prebiotic, so the next move is to ground its required chemical functions in geochemistry. Although documenting this connection requires many individual experimental results, the underlying idea is simple: the laboratory reagents and transformations are ordinary chemistry, and the required chemical functions can be connected through physically realizable sulfur, amino-acid, thioester, and peptide reaction space back to geochemical matter. This gives a geochemically grounded route to a Matsuo-type QAG/self. Variation and selection among such organizations can then provide the upward approach to L*. From the opposite direction, interacting peptide systems derived from modern proteins such as GB1 provide the downward approach toward L*. The proposed solution therefore has three parts: reach L* from modern proteins, reach L* independently from geochemical self-organization, and use the percolating viable network within L* to provide the evolutionary connection between the two entry points.

Note that this third path---the horizontal evolutionary path through L* connecting the QAG-side and modern-protein-side entry points---is not constructed in this paper. Its explicit construction, through viable neighboring peptide organizations connecting these two regimes, is a separate problem and could be a natural subject for future work.

\hypertarget{the-configuration-space-problem}{%
\section*{1. The configuration-space problem}\label{the-configuration-space-problem}}

For a peptide of length N built from 20 amino acids, the raw sequence count 20\^{}N grows too quickly for undirected sampling to be a plausible route to a specified modern protein. But evolution does not need to sample all sequences, and the first selectable peptide organization did not need to be a modern protein. The useful question is whether physically realizable viable configurations can enter a sufficiently connected network for sustained evolutionary traversal.

Let $\Omega$ denote peptide/protein configuration space. A node represents a physically realizable molecular or multicomponent peptide configuration, and an edge represents an accessible change - for example substitution, length change, recruitment or loss of a component, duplication, association, or fusion. We use L* for the peptide/protein configuration-space layer sought where the modern-protein and prebiotic peptide-self-organization approaches converge. Evolutionary manageability is a necessary property of L*: viable configurations within it must belong to a sufficiently connected network for sustained evolutionary traversal.

L* = the sought convergence layer of physically realizable peptide/protein configuration space, with sufficient viable connectivity for sustained evolutionary traversal.

This makes connectivity central. A functional peptide or peptide system that sits on an isolated island is not enough. A system embedded in a connected viable network can provide an evolutionary entrance. The problem is therefore a percolation problem in configuration space: under what peptide-system organization does the accessible viable network cease to consist mainly of disconnected islands and acquire a sufficiently large connected component?

\hypertarget{two-approaches-to-l}{%
\section*{2. Two approaches to L*}\label{two-approaches-to-l}}

The project approaches L* from both sides. Neither route is a discarded alternative; they provide complementary probes of the same evolutionary-manageability problem.

MODERN PROTEINS

$\downarrow$

GB1 (56 aa)

$\downarrow$

viable interacting GB1-derived peptide systems

$\downarrow$

L*

$\uparrow$

functional/evolving peptide networks

$\uparrow$

geochemically grounded QAG / self

$\uparrow$

geochemical matter

The modern-side route uses GB1-derived peptide systems as controlled seeds for probing viable configuration-space connectivity under sequence and compositional variation. The upward route asks how geochemical matter can traverse natural reaction space into a persistent peptide self and how variation among such reproducing organizations can explore peptide configuration space. Agreement or overlap between the connectivity regimes reached from the two directions would constrain L* far more strongly than either approach alone.

\hypertarget{gb1-as-a-modern-side-experimental-probe}{%
\section*{3. GB1 as a modern-side experimental probe}\label{gb1-as-a-modern-side-experimental-probe}}

The B1 domain of streptococcal protein G is a compact 56-residue $\beta$-grasp domain. Kobayashi et al. tested fragment complementation and found that residues 1--40 and 41--56 form a stable 1:1 complex with Kd $\approx$ 9 $\times$ 10\^{}-6 M; NMR showed that the complex resembles the native domain {[}14{]}. Thus a modern protein-like fold can be distributed across independently supplied interacting peptide components rather than requiring one covalent 56-residue chain.

GB1(1--56) -\textgreater{} {[}1--40{]} + {[}41--56{]} -\textgreater{} native-like complex

This is a hard experimental anchor, but not by itself a localization of L*. A fragmentation is relevant to the modern-side construction only if the resulting peptides interact so that, collectively, they preserve the relevant organized protein-like state; cleavage into arbitrary noninteracting peptides does not provide such a system. The demonstrated GB1 system can therefore serve as a controlled seed for probing the viable configurations around it.

Further fragmentation can generate additional GB1-derived interacting peptide seed systems for such connectivity tests. The previously tested split 1--20 + 21--40 + 41--56 did not establish a successful three-fragment system, while computational structural analysis identified nearby candidate partitions 1--19 + 20--40 + 41--56 and 1--21 + 22--40 + 41--56 as wet-chemistry targets. These candidates are predictions, not experimental results. Successful complementation would be useful because it would provide additional independently variable interacting peptide components; it would not by itself imply that the system is closer to L*.

The modern-side GB1 program therefore becomes:

GB1 -\textgreater{} demonstrated interacting 40+16 peptide system -\textgreater{} candidate interacting multipeptide seed systems -\textgreater{} sequence/composition variation -\textgreater{} viable-network connectivity mapping -\textgreater{} L*

The point is not that GB1 itself was ancestral, nor that progressive cleavage moves monotonically toward L*. GB1 is a calibrated modern probe: experimentally viable interacting GB1-derived peptide systems provide starting nodes from which the surrounding configuration space can be varied and its connectivity measured.

\hypertarget{ancient-peptide-modules-strengthen-the-modern-side-approach}{%
\section*{4. Ancient peptide modules strengthen the modern-side approach}\label{ancient-peptide-modules-strengthen-the-modern-side-approach}}

Alva, Söding and Lupas identified a vocabulary of recurrent peptide fragments shared across otherwise unrelated protein folds and proposed that modern domains arose by fusion and accretion of ancestral peptides {[}15{]}. Their candidate fragments span roughly 9--38 residues, with a median near 24 residues. This independently establishes peptide-scale structural material as a plausible substrate for addressing the needle-in-a-haystack problem of protein sequence space, without defining an ordering toward L* by fragment size.

This does not identify L* by itself. It supplies a second modern-side constraint on the needle-in-a-haystack problem: peptide-scale modules provide candidate systems in which viable-network connectivity can be investigated.

modern domains -\textgreater{} recurrent ancient-like peptide modules -\textgreater{} variation/connectivity mapping -\textgreater{} L*

\hypertarget{deacons-autogenic-criterion-qag-ag-oq-and-selfhood}{%
\section*{5. Deacon's autogenic criterion: QAG, AG, OQ, and selfhood}\label{deacons-autogenic-criterion-qag-ag-oq-and-selfhood}}

The upward ladder begins with Deacon's autogenic architecture. Deacon's 2006 autocell model couples two self-organizing processes reciprocally: an autocatalytic component-production process generates material required for a self-assembling enclosure, while enclosure keeps catalytic components in productive proximity. After disruption, released components can restart production and reassemble the organization. Self-reconstitution and self-reproduction can therefore arise from reciprocal linkage without a separate self-replicating molecular template {[}12{]}. In Incomplete Nature Deacon develops this reciprocal constraint logic as part of his broader account of autogenesis and teleodynamic organization {[}13{]}.

We use this Deacon architecture as the conceptual criterion for QAG. A quasi-autogenic organization (QAG) is a non-minimal chemical organization in which reciprocally dependent processes maintain and reproduce the conditions required by the organization {[}12,13{]}. An autogen (AG) is the elementary or minimal case. AG is not more of a self than QAG; the distinction is minimality, not degree of selfhood.

We use organizational QAG (OQ) to denote a non-prebiotic implementation of a QAG: a QAG whose demonstrated components or production chemistry are not yet established as prebiotic.

An OQ is therefore a full QAG and a self. The label concerns the origin or implementation of its chemistry, not the degree of organizational closure or selfhood.

Template-based or sequence-specific heredity is not a prerequisite for QAG selfhood or for the onset of evolution. Deacon's autocell argument shows that reciprocally organized systems can reproduce with variation and undergo a limited form of evolution without self-replicating template molecules {[}12{]}.

\hypertarget{matsuo-supplies-the-qag}{%
\section*{6. Matsuo supplies the QAG}\label{matsuo-supplies-the-qag}}

Matsuo and Kurihara demonstrated a peptide coacervate system in which a synthesized disulfide precursor (Mpre) is reduced by dithiothreitol (DTT) to a reactive cysteine-thioester monomer (M). The monomer spontaneously oligomerizes into a heterogeneous short-peptide population, denoted P, with a reported mean polymerization degree of about 4.1 {[}1{]}.

Mpre -\textgreater{} M -\textgreater{} P

The peptide population P undergoes liquid-liquid phase separation into coacervate droplets, denoted Q. Peptide production continues in association with the droplets and exhibits the concentration-induced autocatalytic behavior described by the authors. Periodic feeding and physical extrusion yield repeated growth-and-division cycles {[}1{]}.

P -\textgreater{} Q -\textgreater{} enhanced production/incorporation of P

The organizational core is therefore:

P \textless-\textgreater{} Q

P supports the compartment; the compartment supports continued production and incorporation of P. This is the same organizational topology Deacon identifies as fundamental: component production and proximity-maintaining self-assembly are reciprocally linked so that the organization can reconstitute and reproduce {[}12,13{]}. The reproduced object in Matsuo is the peptide/coacervate organization, not a specified molecular sequence. Matsuo and Kurihara do not use Deacon's terminology; the classification is ours. Under the Deacon-derived QAG criterion, Matsuo is already a self {[}1,12,13{]}.

Matsuo's published laboratory implementation is therefore an OQ. The organizational closure itself is already demonstrated; the separate question is whether its feed chemistry can be connected downward through physically realizable natural reaction space to geochemical matter.

The geochemical argument below addresses that connection without requiring Nature to reproduce every laboratory reagent or to carry out the entire history in one vessel.

\hypertarget{qag-accessibility-does-not-require-minimality}{%
\section*{7. QAG accessibility does not require minimality}\label{qag-accessibility-does-not-require-minimality}}

Minimality is not a requirement for the present argument. A QAG is already a self-preserving organization, and its functions may be distributed across several components and physical processes. The relevant question is whether such an organization is reachable from geochemical matter and can enter a sufficiently connected peptide configuration space.

The relevant trajectory is:

geochemical matter -\textgreater{} Matsuo-type QAG/self -\textgreater{} variation and selection -\textgreater{} functional peptide networks -\textgreater{} L*

No transition to a minimal autogen is required.

\hypertarget{from-geochemical-matter-to-matsuo-compatible-feed-chemistry}{%
\section*{8. From geochemical matter to Matsuo-compatible feed chemistry}\label{from-geochemical-matter-to-matsuo-compatible-feed-chemistry}}

The bottom of the upward ladder is geochemical matter, not a stipulated 'first peptide' or 'first QAG.' The unknown intermediate chemical history may contain different environments and reaction modules. The relevant existence question is whether a continuous sequence of physically realizable transformations connects geochemical materials to the feed chemistry of the demonstrated Matsuo organization.

geochemical matter \textasciitilde\textgreater{} C/N/S chemistry \textasciitilde\textgreater{} activated amino-acid/thioester chemistry \textasciitilde\textgreater{} Matsuo-compatible feed -\textgreater{} P \textless-\textgreater{} Q

The symbol \textasciitilde\textgreater{} denotes a reaction-space path that may traverse different natural environments. One-pot compatibility is not required upstream; product survival and transfer between successive reaction modules are sufficient. Matsuo's exact laboratory synthesis is therefore not privileged, but neither may its feed simply be assumed: the chemical function performed by each laboratory step must be supplied or bypassed by a supported natural route.

\hypertarget{geochemical-grounding-of-the-feed}{%
\section*{9. Geochemical grounding of the feed}\label{geochemical-grounding-of-the-feed}}

\hypertarget{carbonyl-sulfide-cysteine-and-sulfur-activation}{%
\subsection*{9.1 Carbonyl sulfide, cysteine and sulfur activation}\label{carbonyl-sulfide-cysteine-and-sulfur-activation}}

Leman, Orgel and Ghadiri showed that carbonyl sulfide (COS), a simple sulfur-bearing molecule associated with volcanic chemistry, promotes peptide-bond formation from amino acids in aqueous conditions {[}3{]}. Leman and Ghadiri subsequently demonstrated formation of alpha-amino thioacids in water from alpha-amino acids using COS with sulfide, with N-carboxyanhydride-type activation implicated in the reaction {[}4{]}. Foden et al. demonstrated prebiotically motivated routes into cysteine-containing peptide chemistry {[}2{]}. These results connect simple C/N/S chemistry to activated sulfur/acyl chemistry relevant to the Matsuo feed.

geochemical C/N/S matter \textasciitilde\textgreater{} amino-acid/Cys chemistry + COS/sulfide -\textgreater{} activated sulfur/acyl chemistry

\hypertarget{natural-thiols-and-thioester-reaction-space}{%
\subsection*{9.2 Natural thiols and thioester reaction space}\label{natural-thiols-and-thioester-reaction-space}}

Matsuo uses benzyl mercaptan-derived thioester chemistry as a laboratory implementation and explicitly discusses hydrophobic alkylthiols as prebiotic candidates {[}1{]}. Independent geochemical experiments have demonstrated abiotic methanethiol production from CO2, H2 and H2S on MoS2 under hydrothermal conditions {[}16{]}. Thus the thiol requirement can be connected downward to geochemical carbon, hydrogen and sulfur chemistry rather than to benzyl mercaptan specifically.

CO2 + H2 + H2S + mineral chemistry -\textgreater{} simple thiols

Thioester chemistry itself also has several experimentally supported entrances. Frenkel-Pinter et al. showed that mercaptoacids and amino acids form thioester intermediates and proto-peptides during wet-dry chemistry {[}5{]}. Weber demonstrated in water that activated amino-acid/N-carboxyanhydride chemistry can be intercepted by a thiol to produce peptide thioesters {[}17{]}. The latter experiment used CDI for activation and therefore is not itself a prebiotic synthesis; its relevance here is the experimentally demonstrated NCA-to-thioester reaction topology, while COS chemistry supplies an independent prebiotic route into NCA-type amino-acid activation {[}3,4{]}.

\hypertarget{the-dry-reaction-space-splice}{%
\subsection*{9.3 The dry reaction-space splice}\label{the-dry-reaction-space-splice}}

Taken together, the literature supports a reaction-space construction of the form: geochemical matter \textasciitilde\textgreater{} Cys/amino-acid chemistry + COS/sulfide + natural thiols \textasciitilde\textgreater{} activated amino-acid sulfur chemistry \textasciitilde\textgreater{} Matsuo-compatible thioester feed \textasciitilde\textgreater{} peptides. The exact unprotected cysteine-thioester sequence has not been demonstrated end-to-end under one prebiotic protocol; the construction therefore claims chemical connectivity, not a completed one-pot reconstruction. The remaining uncertainty at the splice is reaction selectivity and kinetics, not the absence of a chemically grounded precursor class.

\hypertarget{peptide-material---q-short-peptide-coacervation}{%
\section*{10. Peptide material -\textgreater{} Q: short peptide coacervation}\label{peptide-material---q-short-peptide-coacervation}}

Prebiotic access to the aromatic amino-acid component is independently supported by Friedmann and Miller, who demonstrated phenylalanine and tyrosine synthesis under simulated primitive-Earth conditions {[}6{]}.

Han et al. systematically examined short Phe-rich tripeptides and found that FGF, FFG, FAF and FFA form liquid coacervate droplets under the reported conditions, whereas closely related sequences can instead aggregate {[}7{]}. This establishes liquid coacervation at only three residues and makes FFG and FGF concrete model members of the target product space.

FGF -\textgreater{} Q FFG -\textgreater{} Q

A primitive peptide synthesis is unlikely to produce a purified single sequence. That does not invalidate the route. Cao et al. showed that binary mixtures of related short aromatic peptides can stabilize liquid coacervates and suppress conversion into more ordered states {[}8{]}. A heterogeneous peptide population can therefore be a feature of the QAG rather than noise to be eliminated.

P* = a heterogeneous short Q-forming peptide population; Matsuo P is itself a heterogeneous peptide population

\hypertarget{q---peptide-production-independent-precedent}{%
\section*{11. Q -\textgreater{} peptide production: independent precedent}\label{q---peptide-production-independent-precedent}}

Wang et al. demonstrated selective amide-bond formation in redox-active coacervate protocells using sulfur-mediated chemistry {[}9{]}. The coacervate environment concentrates reactants and enhances peptide ligation relative to bulk solution. Phenylalanine was tested, but its comparatively strong interaction with the coacervate matrix reduced its reactivity relative to several other amino acids. The important result for the present argument is therefore the general coacervate enhancement of peptide ligation, not preferential phenylalanine incorporation. This establishes the mechanism needed for the return arrow:

Q -\textgreater{} enhanced peptide production

The Wang coacervate is not formed by the same peptide population being synthesized, so it is not an independent reconstruction of Matsuo closure. Its importance is narrower and positive: coacervate-enhanced sulfur-mediated peptide synthesis is experimentally real and therefore supports the physical plausibility of the return mechanism already demonstrated organizationally by Matsuo.

\hypertarget{compatibility-and-intersection-of-condition-spaces}{%
\section*{12. Compatibility and intersection of condition spaces}\label{compatibility-and-intersection-of-condition-spaces}}

A reaction-space path is not useful if its terminal chemistry and the QAG require mutually impossible environments. We therefore separate upstream compartmentalization from the stricter compatibility requirement at the Matsuo interface.

\begin{quote}
Upstream geochemical reactions may occupy different spatial or temporal compartments: hydrothermal synthesis, mineral-surface chemistry, transport, drying, rehydration and pond chemistry need not occur simultaneously.

Let C\_feed denote the condition space supporting the terminal sulfur-activated feed chemistry and C\_Matsuo the condition space supporting M -\textgreater{} P -\textgreater{} Q -\textgreater{} continued P production. The relevant interface test is whether C\_feed intersects C\_Matsuo.

C\_feed intersect C\_Matsuo != empty

The demonstrated chemistries provide such an overlap at the level of the major environmental variables: aqueous conditions, moderate temperature, and neutral-to-mildly alkaline pH are occupied by COS-mediated amino-acid chemistry and lie within the broad aqueous pH regime in which Matsuo reports droplet formation {[}1,3,4{]}.

Accordingly, there is no identified water/temperature/pH barrier forcing terminal feed production and the Matsuo organization into mutually exclusive environments. The exact product distribution and kinetics of the proposed cysteine/thiol splice remain experimentally testable, but this is not a missing reaction-space bridge.
\end{quote}

The resulting construction is:

compartmentalized geochemical chemistry \textasciitilde\textgreater{} common aqueous interface -\textgreater{} Matsuo-compatible feed -\textgreater{} P \textless-\textgreater{} Q

This is the dry geogrounding result. It does not assert that the complete path has been reconstructed experimentally in one vessel, nor does the existence argument require such a reconstruction.

\hypertarget{from-geochemical-matter-toward-l}{%
\section*{13. From geochemical matter toward L*}\label{from-geochemical-matter-toward-l}}

The upward ladder is therefore no longer a speculative jump from chemistry directly to proteins. It contains experimentally anchored stages:

L*

$\uparrow$

evolving functional peptide configurations

$\uparrow$

selection among persistent QAG variants

$\uparrow$

variation among reproducing QAGs

$\uparrow$

geochemically grounded QAG / self

$\uparrow$

short peptide population P*

$\uparrow$

geochemical matter

Before a QAG exists, successful chemical combinations are retained only by environmental recurrence. The geochemical route may therefore sample through multiple reaction compartments before entering Matsuo-compatible chemistry. After a QAG exists, the organization itself can reproduce with variation. As Deacon emphasizes, such variant organizations can differentially propagate even without self-replicating template molecules {[}12{]}. This provides the transition from geochemical sampling to selection among persistent chemical organizations.

\subsection*{13.1 Sequence heredity as a later strengthening}

Sequence-specific heredity is not required for QAG selfhood or for the initial onset of selection among reproducing QAG organizations. It can nevertheless arise later and may strengthen the retention of particular peptide innovations.

One candidate sequence-heredity mechanism uses a replicating peptide sequence R coupled to the QAG compartment Q:

R -\textgreater{} R R -\textgreater{} Q Q -\textgreater{} R

Here R -\textgreater{} R denotes persistence of a peptide sequence by replication; R -\textgreater{} Q means that the sequence affects persistence or reproduction of the QAG; and Q -\textgreater{} R means that the QAG environment supports continued replication of that sequence. This is a possible later mechanism of sequence-specific heredity, not a prerequisite for QAG selfhood or for the beginning of evolution.

Nanda et al. demonstrated a nonenzymatic peptide replication network based on thioester ligation and beta-sheet templating, providing an experimental peptide-replication precedent {[}10{]}. Baba et al. showed that specific short peptide sequences can alter fatty-acid vesicle growth, providing a separate precedent for sequence-dependent effects on compartment performance {[}11{]}. These are different systems and do not yet demonstrate sequence heredity coupled to the same QAG.

\hypertarget{percolation-and-evolutionary-manageability}{%
\section*{14. Percolation and evolutionary manageability}\label{percolation-and-evolutionary-manageability}}

The central configuration-space claim is that evolutionary accessibility depends on network connectivity, not on peptide length or any assumed scalar ordering of peptide systems. Viable configurations may exist yet remain sparse and isolated. Other peptide-system organizations may support many mutationally or compositionally adjacent viable states and thereby form a large connected component.

For a candidate peptide-system ensemble E, let G(E) be the graph of viable configurations accessible within that ensemble. A useful order parameter is the relative size of its largest connected component:

S(E) = \textbar largest connected component of G(E)\textbar{} / \textbar G(E)\textbar{}

The transition need not be mathematically sharp in a finite chemical system. Operationally, a candidate layer is evolutionarily manageable when S(E) is large enough that successful configurations have many viable routes to other successful configurations rather than being isolated endpoints. Evolutionary manageability is therefore a necessary criterion for identifying L*.

This gives the two approaches a common test. On the modern side, GB1-derived and ancient-like peptide systems provide seed configurations around which viable sequence/composition connectivity can be mapped. On the geochemical side, the question is whether variation among QAG-supported peptide organizations can enter a comparably traversable viable network. The two approaches are directed toward identifying the same L*.

\hypertarget{convergence-of-the-two-approaches}{%
\section*{15. Convergence of the two approaches}\label{convergence-of-the-two-approaches}}

MODERN PROTEINS

$\downarrow$

GB1

$\downarrow$

viable interacting peptide seed systems

$\downarrow$

sequence/composition variation

$\downarrow$

viable-network connectivity mapping

$\downarrow$

L*

$\uparrow$

functional peptide networks

$\uparrow$

reproducing/selected QAG variants

$\uparrow$

prebiotic QAG

$\uparrow$

peptide P*

$\uparrow$

geochemical matter

No numerical scale for L* is assumed. The GB1 complementation result and Alva's recurrent ancient fragments provide experimentally and evolutionarily motivated peptide-scale seed systems on the modern side {[}14,15{]}. The Matsuo/prebiotic route shows that selfhood can begin as an organization of very short heterogeneous peptides {[}1,7{]}. The remaining question is whether variation from these respective starting regimes reaches a sufficiently connected viable network for sustained evolutionary traversal.

This is why both approaches remain in the same project. The upward route supplies a plausible entrance into selfhood and selection; the modern-side route supplies experimentally tractable peptide systems from which surrounding viable configuration-space connectivity can be measured; percolation supplies the common criterion for evolutionary manageability.

\hypertarget{experimental-program}{%
\section*{16. Experimental program}\label{experimental-program}}

The work separates into a modern-side connectivity-mapping program, an upward evolutionary-connectivity program, an optional later heredity branch, and experimental tests that can strengthen the geochemical-to-QAG construction.

A. Geochemical-QAG validation. Test the strongest dry splice directly: generate Matsuo-compatible thioester feed from geochemically grounded Cys/amino-acid, COS/sulfide and natural-thiol chemistry within the overlapping aqueous condition space, then test entry into the Matsuo P \textless-\textgreater{} Q regime. This is experimental validation of an already constructed reaction-space path, not a prerequisite for the dry connectivity claim.

B. Modern-side connectivity mapping. Use the demonstrated interacting GB1 1-40 + 41-56 system and, if experimentally viable, the candidate interacting three-fragment systems 1-19 + 20-40 + 41-56 and 1-21 + 22-40 + 41-56 as seed configurations. Vary component sequences and compositions and map which neighboring combinations remain viable, rather than treating successful fragmentation itself as movement toward L*.

Successful three-fragment complementation would therefore add useful interacting peptide seed systems and independently variable components; evidence relevant to L* would come from the connectivity of viable variants around those seeds.

C. Later sequence-heredity branch. If stronger sequence-specific inheritance is required, test whether a replicating peptide sequence can both affect QAG performance and receive replication support from the same QAG. This is not required to establish the prebiotic QAG or the onset of selection.

In parallel, computational graph construction can estimate connectivity across peptide substitutions, fragment boundaries, association states, duplication and fusion events. The experimental and computational tasks then meet at the same object: the connected viable network.

\hypertarget{conclusion}{%
\section*{17. Conclusion}\label{conclusion}}

The project now separates experimental demonstrations from the reaction-space construction connecting them. Deacon supplies the conceptual criterion: reciprocal coupling of component production and self-assembling proximity maintenance can generate a self-reconstituting, self-reproducing organization without prior template replication {[}12,13{]}. Matsuo supplies an experimental peptide/coacervate realization of that topology {[}1{]}. The published laboratory implementation is therefore an organizational QAG (OQ), while the geochemical analysis below connects its feed requirements downward to geochemical matter through physically realizable reaction space.

What is established:

1. Geochemical and prebiotic chemistry provides experimentally demonstrated entrances into cysteine/amino-acid, COS/sulfide, natural-thiol, thioester and short-peptide reaction space {[}2-5,16,17{]}.

2. Very short phenylalanine-rich peptides can form liquid coacervate droplets, and heterogeneous short-peptide mixtures can stabilize liquid condensates {[}7,8{]}.

3. Coacervate environments can enhance sulfur-mediated peptide ligation, establishing a physically plausible return mechanism from compartment formation to increased peptide production {[}9{]}.

4. Matsuo demonstrates the complete organizational topology P \textless-\textgreater{} Q with laboratory feed: peptide production generates the coacervate phase, and the coacervate supports continued peptide production and incorporation {[}1{]}.

5. From the modern-protein side, GB1 demonstrates that a 56-residue protein-like fold can be distributed across independently supplied interacting 40- and 16-residue peptide fragments {[}14{]}. This provides an experimentally controlled interacting peptide system from which neighboring viable configurations can be probed. Alva et al. identify recurrent ancient peptide fragments in roughly the 9-38-residue range, providing additional peptide-scale material relevant to the modern-side connectivity question {[}15{]}.

6. Separate experiments establish peptide replication and sequence-dependent effects on compartment growth {[}10,11{]}. They support the plausibility of later sequence-specific heredity but do not yet constitute a sequence-heritable QAG.

What remains experimentally open:

1. The geochemical-to-Matsuo reaction-space construction has not yet been executed end-to-end experimentally. In particular, the strongest terminal splice into Matsuo-compatible cysteine-thioester feed remains a direct selectivity/kinetics test. This is an experimental validation question rather than a missing dry bridge.

2. On the modern side, viable-network connectivity around GB1-derived interacting peptide systems has not yet been mapped. The proposed three-fragment partitions remain useful wet-chemistry candidates as additional interacting peptide seed systems, but successful complementation alone would not locate L*.

The origin problem is therefore recast through percolation as the problem of moving from geochemical matter, through a geochemically grounded Matsuo-type QAG/self, into a sufficiently connected peptide network from which sustained evolution toward protein-like organization becomes possible.

\hypertarget{references}{%
\section*{References}\label{references}}

{[}1{]} Matsuo, M. \& Kurihara, K. (2021). Proliferating coacervate droplets as the missing link between chemistry and biology in the origins of life. Nature Communications 12, 5487. doi:10.1038/s41467-021-25530-6.

{[}2{]} Foden, C. S. et al. (2020). Prebiotic synthesis of cysteine peptides that catalyze peptide ligation in neutral water. Science 370, 865--869. doi:10.1126/science.abd5680.

{[}3{]} Leman, L., Orgel, L. \& Ghadiri, M. R. (2004). Carbonyl Sulfide-Mediated Prebiotic Formation of Peptides. Science 306, 283--286. doi:10.1126/science.1102722.

{[}4{]} Leman, L. J. \& Ghadiri, M. R. (2016, published online 9 November; Synlett 2017, 28(01), 68-72). Potentially Prebiotic Synthesis of alpha-Amino Thioacids in Water. doi:10.1055/s-0036-1589410.

{[}5{]} Frenkel-Pinter, M. et al. (2022). Thioesters provide a plausible prebiotic path to proto-peptides. Nature Communications 13, 2569. doi:10.1038/s41467-022-30191-0.

{[}6{]} Friedmann, N. \& Miller, S. L. (1969). Phenylalanine and tyrosine synthesis under primitive earth conditions. Science 166, 766--767. doi:10.1126/science.166.3906.766.

{[}7{]} Han, H. et al. (2026). Sequence-Modulated Active Tripeptide Condensates for Tandem Catalysis. Angewandte Chemie International Edition 65, e17620. doi:10.1002/anie.202517620.

{[}8{]} Cao, S. et al. (2025). Binary peptide coacervates as an active model for biomolecular condensates. Nature Communications 16, 2407. doi:10.1038/s41467-025-57772-z.

{[}9{]} Wang, J., Abbas, M., Wang, J. \& Spruijt, E. (2023). Selective amide bond formation in redox-active coacervate protocells. Nature Communications 14, 8492. doi:10.1038/s41467-023-44284-x.

{[}10{]} Nanda, J. et al. (2017). Emergence of native peptide sequences in prebiotic replication networks. Nature Communications 8, 434. doi:10.1038/s41467-017-00463-1.

{[}11{]} Baba, A. et al. (2026). Growth of fatty acid vesicles coupled with amino acid sequences of peptides toward evolvable protocells. Communications Chemistry 9, 234. doi:10.1038/s42004-026-02043-1.

{[}12{]} Deacon, T. W. (2006). Reciprocal Linkage between Self-organizing Processes is Sufficient for Self-reproduction and Evolvability. Biological Theory 1(2), 136-149. doi:10.1162/biot.2006.1.2.136.

{[}13{]} Deacon, T. W. (2011). Incomplete Nature: How Mind Emerged from Matter. W. W. Norton \& Company.

{[}14{]} Kobayashi, N., Honda, S., Yoshii, H., Uedaira, H. \& Munekata, E. (1995). Complement assembly of two fragments of the streptococcal protein G B1 domain in aqueous solution. FEBS Letters 366, 99--103. doi:10.1016/0014-5793(95)00503-2.

{[}15{]} Alva, V., Söding, J. \& Lupas, A. N. (2015). A vocabulary of ancient peptides at the origin of folded proteins. eLife 4:e09410. doi:10.7554/eLife.09410.

{[}16{]} Kitadai, N., Shibuya, T., Ueda, H., Tasumi, E., Okada, S. \& Takai, K. (2024). Supercritical carbon dioxide likely served as a prebiotic source of methanethiol in primordial ocean hydrothermal systems. Communications Earth \& Environment 5, 539. doi:10.1038/s43247-024-01689-w.

{[}17{]} Weber, A. L. (2005). Aqueous synthesis of peptide thioesters from amino acids and a thiol using 1,1'-carbonyldiimidazole. Origins of Life and Evolution of Biospheres 35, 421-427. doi:10.1007/s11084-005-4070-0.

\end{document}
